% deckblatt.tex, 2019/03/15 
\documentclass[a4paper,12pt]{report}  
%Inhaltsverzeichnis in Kapitel und Sektionen aufgeteilt  \chapter{title} > \section{title} > subsection{title} \subsubsection{title}

% Pakete und Paket-Configs
\usepackage{times} % Times Roman als Standardschrift
\usepackage[ngerman]{babel} % neue deutsche Rechtschreibung und Trennung
\usepackage{fancyhdr} % spezielle Kopfzeilen
\usepackage[latin1]{inputenc} % Umlaute ��� auch normal benutzen und nicht maskieren
\usepackage[babel, german=quotes]{csquotes}
\usepackage{subfigure} % Figures divided into subfigures.
\usepackage{ifthen} % Ermöglicht ifthenelse und whiledo
\usepackage{amsfonts} % Extra mathematical symbols
\usepackage[rflt]{floatflt} % verbessertes floatfig, als um figure's fliesende texte
\usepackage[T1]{fontenc} % ?? aber notwendig für korrekte PDF-Metadaten
%\usepackage{longtable} % Support for tables longer than a page.
%\usepackage{a4wide} % Increases width of printed area of an a4 page.
%\usepackage{alltt} % verbatim environment except that \ and braces have their usual meanings.
\usepackage{listings} % Typeset source code listings using LaTeX.
\usepackage{moreverb} % bessere verbatim-umgebungen
\usepackage{graphicx}
\usepackage{pdfpages}
\usepackage{xcolor} % Farben Definierbar \definecolor{fhorange}{RGB}{255,153,0}
\usepackage{biblatex} % F�r das Erstellen eines Literaturverzeichnisses
\addbibresource{literatur.bib}  % f�r das erstellen des Literaturverzeichnisses empfielt sich die Software Mendeley, diese erm�glicht einfaches hinzuf�gen B�chern und das exportieren in eine .bib Datei 



\pagestyle{headings} % Kopf- und Fusszeilen
\pagenumbering{Roman} % Nummerierung der Seiten
\newcommand{\maximagewidth}{15cm} % maximal m�gliche Bildbreite

\setlength{\parindent}{0cm} % Einr�ckung am Abstzanfang
\setlength{\parskip}{5pt plus 2pt minus 1pt} % Abstand der Abs�tze zueinander
\frenchspacing % Kein Zusatzzwischenraum nach Satzzeichen

\setcounter{secnumdepth}{3} % Z�hlung bis paragraph 1.1.1
\setcounter{tocdepth}{3} % Inhaltsverzeichnis bis paragraph 1.1.1

%\makeglossary % Schreibe ein Glossar-File

\title{Schaffung einer fusionierten Datenbasis zur Immobilienbewertung in Microsoft Azure}
\author{Andreas Hofmann}
\date{\today}

\definecolor{fhorange}{RGB}{255,153,0}
\begin{document}

%%% Titelseite %%%
\begin{titlepage}
\thispagestyle{empty}
\centerline{
\begin{minipage}{15cm}
	\flushright{
		\includegraphics[width=4cm]{fhlogo.png} \hspace{1.5cm} \ \\
		\textbf{\large Hochschule \hspace{3.5cm} ~ \\ Augsburg } \large University of \hspace{1cm} ~ \\ Applied Sciences \hspace{0.2cm} ~
		\\~
		}
	\flushleft{
		\textbf{\large \textcolor{fhorange}{Bachelorthesis}} \hspace{7.7cm} \textcolor{fhorange}{\large Fakult�t f�r}
		\vspace{0.1cm}
		\\~  \hspace{10.75cm} \textcolor{fhorange}{\large Informatik} 
	}
	%\vspace{1cm}
	\flushleft{ \large Studienrichtung \\ 
	%	\vspace{0.2cm}
		Informatik
	}
	\vspace{0.1cm}
	\flushleft{
	\textcolor{fhorange}{\large \textbf{Andreas Hofmann \\
	Schaffung einer fusionierten Datenbasis zur \\Immobilienbewertung in Microsoft Azure}}
}
	\vspace{0.2cm}
	\flushleft{ \large Erstpr�fer: Dipl.-Inf. (FH), Dipl.-De Erich Seifert 
	\\ Zweitpr�fer: Prof. Dr. Rolf Winter\\ 
		\vspace{0.1cm}
		Abgabe der Arbeit am: 20.08.2021			
	}
	\flushleft{ %In Kooperation mit Firma: \\ Hsa-Digit \\ Betreuer: Hans Mustermann \\ 
		%\vspace{0.5cm}
		%\includegraphics[width=3cm]{hsa_digit.png} \hspace{5cm} 
		\tiny{\\ \textcolor{gray}{
			\hspace{11.3cm} Hochschule f�r angewandte \\
			\hspace{11.3cm} Wissenschaften Augsburg \\
			\hspace{11.3cm} University of Applied Sciences \\ ~ \\
			\hspace{11.3cm} An der Hochschule 1 \\
			\hspace{11.3cm} D-86161 Augsburg \\ ~ \\
			\hspace{11.3cm} Telefon +49 821 55 86-0 \\
			\hspace{11.3cm} Fax +49 821 55 86-3222 \\
			\hspace{11.3cm} www.hs-augsburg.de \\
			\hspace{11.3cm} info@hs-augsburg.de \\} ~ \\
			\hspace{11.3cm} Fakult�t f�r Informatik \\
			\hspace{11.3cm} Telefon +49 821 5586-3450 \\
			\hspace{11.3cm} Fax ~~~~~ +49 821 5586-3499 \\ ~ \\
			\hspace{11.3cm} Verfasser der Bachelorthesis: \\
			\hspace{11.3cm} Max Mustermann \\
			\hspace{11.3cm} Musterstr. 123abc \\
			\hspace{11.3cm} 12345 Musterhausen \\
			\hspace{11.3cm} Telefon +49 123 456 789 \\
			\hspace{11.3cm} max.musterman@musterhost.de \\
	}}
% Nicht �ber die Fehlermeldung There�s no line here to end. \end... liegt daran LaTex es nicht mag wenn Leerzeilen manuell eingef�gt werden
\end{minipage}}\par
\end{titlepage} 


% Zusammenfassung
\thispagestyle{empty}
\section*{Abstract}
...

\newpage
\tableofcontents  
% Wenn das Inhaltsverzeichnis aktualisiert werden muss 2 mal kompilieren, erst dann wird die Aktualisierung angezeigt
\newpage
\pagenumbering{arabic}


%% Include alle Sub-Files
\chapter{Einleitung}
	Lorem ipsum dolor sit amet, consetetur sadipscing elitr, sed diam nonumy eirmod tempor invidunt ut labore et dolore magna aliquyam erat, sed diam voluptua. At vero eos et accusam et justo duo dolores et ea rebum. Stet clita kasd gubergren, no sea takimata sanctus est Lorem ipsum dolor sit amet. Lorem ipsum dolor sit amet, consetetur sadipscing elitr, sed diam nonumy eirmod tempor invidunt ut labore et dolore magna aliquyam erat, sed diam voluptua. At vero eos et accusam et justo duo dolores et ea rebum. Stet clita kasd gubergren, no sea takimata sanctus est Lorem ipsum dolor sit amet. Lorem ipsum dolor sit amet, consetetur sadipscing elitr, sed diam nonumy eirmod tempor invidunt ut labore et dolore magna aliquyam erat, sed diam voluptua. At vero eos et accusam et justo duo dolores et ea rebum. Stet clita kasd gubergren, no sea takimata sanctus est Lorem ipsum dolor sit amet.   
	
	Duis autem vel eum iriure dolor in hendrerit in vulputate velit esse molestie consequat, vel illum dolore eu feugiat nulla facilisis at vero eros et accumsan et iusto odio dignissim qui blandit praesent luptatum zzril delenit augue duis dolore te feugait nulla facilisi. Lorem ipsum dolor sit amet, consectetuer adipiscing elit, sed diam nonummy nibh euismod tincidunt ut laoreet dolore magna aliquam erat volutpat.   
	

\chapter{Grundlagen}
	Lorem ipsum dolor sit amet, consetetur sadipscing elitr, sed diam nonumy eirmod tempor invidunt ut labore et dolore magna aliquyam erat, sed diam voluptua. At vero eos et accusam et justo duo dolores et ea rebum. Stet clita kasd gubergren, no sea takimata sanctus est Lorem ipsum dolor sit amet. Lorem ipsum dolor sit amet, consetetur sadipscing elitr, sed diam nonumy eirmod tempor invidunt ut labore et dolore magna aliquyam erat, sed diam voluptua. At vero eos et accusam et justo duo dolores et ea rebum. Stet clita kasd gubergren, no sea takimata sanctus est Lorem ipsum dolor sit amet. Lorem ipsum dolor sit amet, consetetur sadipscing elitr, sed diam nonumy eirmod tempor invidunt ut labore et dolore magna aliquyam erat, sed diam voluptua. At vero eos et accusam et justo duo dolores et ea rebum. Stet clita kasd gubergren, no sea takimata sanctus est Lorem ipsum dolor sit amet.   
	
\section{Grundlagen in A}
	Lorem ipsum dolor sit amet, consetetur sadipscing elitr, sed diam nonumy eirmod tempor invidunt ut labore et dolore magna aliquyam erat, sed diam voluptua. At vero eos et accusam et justo duo dolores et ea rebum. Stet clita kasd gubergren, no sea takimata sanctus est Lorem ipsum dolor sit amet. Lorem ipsum dolor sit amet, consetetur sadipscing elitr, sed diam nonumy eirmod tempor invidunt ut labore et dolore magna aliquyam erat, sed diam voluptua. At vero eos et accusam et justo duo dolores et ea rebum. Stet clita kasd gubergren, no sea takimata sanctus est Lorem ipsum dolor sit amet. Lorem ipsum dolor sit amet, consetetur sadipscing elitr, sed diam nonumy eirmod tempor invidunt ut labore et dolore magna aliquyam erat, sed diam voluptua. At vero eos et accusam et justo duo dolores et ea rebum. Stet clita kasd gubergren, no sea takimata sanctus est Lorem ipsum dolor sit amet. 
	
\subsection{Grundlagen in B in A}

	Lorem ipsum dolor sit amet, consetetur sadipscing elitr, sed diam nonumy eirmod tempor invidunt ut labore et dolore magna aliquyam erat, sed diam voluptua. At vero eos et accusam et justo duo dolores et ea rebum. Stet clita kasd gubergren, no sea takimata sanctus est Lorem ipsum dolor sit amet. Lorem ipsum dolor sit amet, consetetur sadipscing elitr, sed diam nonumy eirmod tempor invidunt ut labore et dolore magna aliquyam erat, sed diam voluptua. At vero eos et accusam et justo duo dolores et ea rebum. Stet clita kasd gubergren, no sea takimata sanctus est Lorem ipsum dolor sit amet. Lorem ipsum dolor sit amet, consetetur sadipscing elitr, sed diam nonumy eirmod tempor invidunt ut labore et dolore magna aliquyam erat, sed diam voluptua. At vero eos et accusam et justo duo dolores et ea rebum. Stet clita kasd gubergren, no sea takimata sanctus est Lorem ipsum dolor sit amet. 
	
\subsubsection{Grunlagen in C in B in A}

	Lorem ipsum dolor sit amet, consetetur sadipscing elitr, sed diam nonumy eirmod tempor invidunt ut labore et dolore magna aliquyam erat, sed diam voluptua. At vero eos et accusam et justo duo dolores et ea rebum. Stet clita kasd gubergren, no sea takimata sanctus est Lorem ipsum dolor sit amet. Lorem ipsum dolor sit amet, consetetur sadipscing elitr, sed diam nonumy eirmod tempor invidunt ut labore et dolore magna aliquyam erat, sed diam voluptua. At vero eos et accusam et justo duo dolores et ea rebum. Stet clita kasd gubergren, no sea takimata sanctus est Lorem ipsum dolor sit amet. Lorem ipsum dolor sit amet, consetetur sadipscing elitr, sed diam nonumy eirmod tempor invidunt ut labore et dolore magna aliquyam erat, sed diam voluptua. At vero eos et accusam et justo duo dolores et ea rebum. Stet clita kasd gubergren, no sea takimata sanctus est Lorem ipsum dolor sit amet. 


%\include{4anforderungen}
%\include{3stand}
%\include{5spezifikation}
%\include{6implementation}
%\include{7zusammenfassung}
%\include{8suffix}
% Inhaltsverzeichnis




\end{document}
